\chapter{Exchange-Traded Funds}\label{ch:m1:exchange-traded-funds}

A fund's shares trade at \$50.10 on the exchange while the five hundred
stocks inside it are worth \$50.00 a share. Within a second or two a trading
firm sells the fund, buys the five hundred stocks, and tonight hands those
stocks to the fund in exchange for newly printed fund shares that close its
short. It has earned a few cents a share; the fund has grown; the price is
back in line. Nobody at the fund decided anything. At the end of August 2026
some twenty-four trillion dollars were invested in vehicles kept honest by
this one mechanism. This chapter explains the mechanism, what happens to it
when the stocks inside are closed or cannot be priced, and the funds that
promise three times the index and deliver something else.

\section{A fund with two markets}

\begin{definition}[Exchange-traded fund]\label{def:m1:exchange-traded-funds:etf}
An \emph{exchange-traded fund}\index{exchange-traded fund} (ETF) is an
open-ended investment fund whose shares are listed and trade on exchanges
throughout the day like any stock, and can in addition be created or redeemed
with the fund itself, in large blocks, by designated intermediaries.
\end{definition}

\begin{definition}[Authorised participant, creation unit, creation basket]\label{def:m1:exchange-traded-funds:ap}
An \emph{authorised participant}\index{authorised participant} (AP) is a
\omterm{def:m1:the-sell-side:sell-side}{broker-dealer} that has contracted with an ETF for the right to create and
redeem its shares. It does so in \emph{creation units}\index{creation unit},
blocks of a fixed number of ETF shares, by delivering to the fund (or
receiving from it) the \emph{creation basket}\index{creation basket}: the
list of securities and cash, published daily by the fund, that one creation
unit is exchanged for.
\end{definition}

\begin{definition}[In-kind transfer]\label{def:m1:exchange-traded-funds:inkind}
A creation or redemption is \emph{in kind}\index{in-kind transfer} when the
AP and the fund exchange ETF shares for the securities themselves, not for
cash. The fund then never trades in the market: the cost of turning cash
into a portfolio, and back, is borne by the investors who come and go, not by
those who stay.
\end{definition}

\begin{omfigure}
\begin{tikzpicture}[font=\small,
  box/.style={draw=omDef,rounded corners=2pt,minimum width=2.9cm,minimum height=1.0cm,align=center,fill=omDef!4},
  ext/.style={draw=omExo,rounded corners=2pt,minimum width=2.9cm,minimum height=1.0cm,align=center,fill=omExo!6},
  flow/.style={-{Stealth[length=2mm]},thick}]
\node[ext] (inv) at (0,1.4) {Investors};
\node[box] (ex) at (5.2,1.4) {Exchange\\(secondary market)};
\node[box] (ap) at (10.4,1.4) {Authorised participant\\and market makers};
\node[box] (fund) at (10.4,-1.6) {The fund\\(primary market)};
\node[ext] (und) at (5.2,-1.6) {Market for the\\underlying securities};
\draw[flow,omDef] ([yshift=5pt]inv.east) -- node[above,font=\footnotesize]{trade} ([yshift=5pt]ex.west);
\draw[flow,omDef] ([yshift=5pt]ap.west) -- node[above,font=\footnotesize]{quote} ([yshift=5pt]ex.east);
\draw[flow,omThm] ([xshift=-7pt]ap.south) -- node[left,font=\footnotesize,align=right]{basket in,\\ETF shares out} ([xshift=-7pt]fund.north);
\draw[flow,omThm] ([xshift=7pt]fund.north) -- node[right,font=\footnotesize,align=left]{or the\\reverse} ([xshift=7pt]ap.south);
\draw[flow,omDef] ([yshift=-8pt]ap.west) -- ++(-1.0,0) |- node[pos=0.25,left,font=\footnotesize,align=right]{buys or sells\\the basket} (und.east);
\end{tikzpicture}

\omcaption{The two markets of an ETF. Investors meet only the exchange. The
number of ETF shares in existence changes only on the right-hand side, in
\omterm{def:m1:exchange-traded-funds:ap}{creation units}, at the day's \omterm{def:m1:exchange-traded-funds:nav}{net asset value}, against the
basket.}\label{fig:m1:exchange-traded-funds:two}
\end{omfigure}

\begin{dated}{2026-09}{Size, and the rule that standardised the structure}\label{dat:m1:exchange-traded-funds:size}
Assets in \omterm{def:m1:exchange-traded-funds:etf}{exchange-traded funds} worldwide reached \$24.03 trillion at the end
of August 2026, a record, by an industry consultancy's count. In the United States, Rule 6c-11 under the
Investment Company Act, adopted on 25 September 2019, lets a transparent ETF
operate without an individual exemptive order, provided it publishes its
portfolio daily, and allows \emph{custom baskets} that differ from a pro-rata
slice of the portfolio. Leveraged and inverse ETFs cannot rely on the rule.
\end{dated}

\section{The arbitrage that holds the price}

\begin{definition}[Net asset value, indicative value, premium]\label{def:m1:exchange-traded-funds:nav}
The \emph{net asset value}\index{net asset value} (NAV) of a fund share is
the value of the fund's holdings less its liabilities, divided by the number
of shares, computed once a day at official closing prices. The
\emph{indicative NAV}\index{indicative NAV} is an estimate of the same
quantity recomputed during the day from current prices of the basket. The
\emph{premium}\index{premium} of an ETF is $P/\mathrm{NAV} - 1$, with $P$ its
market price; a negative premium is a \emph{discount}.
\end{definition}

\begin{proposition}[The arbitrage band]\label{prop:m1:exchange-traded-funds:band}
Let $c_B$ and $c_E$ be the costs of trading the basket and the ETF (half the
spread plus impact, as fractions of value), $\phi$ the creation or redemption
fee per unit of value, and $f$ the cost of financing both legs until the
creation settles. An \omterm{def:m1:exchange-traded-funds:ap}{authorised participant} profits from creating when the
\omterm{def:m1:exchange-traded-funds:nav}{premium} exceeds $b = c_B + c_E + \phi + f$ and from redeeming when the
discount exceeds $b$. In a market with competing APs the \omterm{def:m1:exchange-traded-funds:nav}{premium} therefore
stays within $[-b, +b]$.
\end{proposition}
\begin{proof}
At \omterm{def:m1:exchange-traded-funds:nav}{premium} $\pi > 0$ the AP sells ETF shares at $P(1 - c_E)$ and buys the
basket at $\mathrm{NAV}(1 + c_B)$, pays the fee and the financing, and
delivers the basket for ETF shares at NAV. Its profit per unit of value is
$\pi - c_E - c_B - \phi - f$ to first order, positive iff $\pi > b$. The
redemption trade is symmetric.
\end{proof}

\begin{example}[Seven basis points]\label{ex:m1:exchange-traded-funds:band}
For a large-capitalisation equity ETF take $c_B = 4$, $c_E = 1.5$, $\phi = 1$
and $f = 0.5$ basis points: $b = 7$ basis points, three and a half cents on a
\$50 share. For a fund of small or foreign stocks, $c_B$ alone can be thirty
basis points; for a fund of corporate bonds, more. The band is not a defect
of the ETF: it is the cost of trading the underlying, which the ETF investor
avoids so long as the price stays inside it
(\cref{fig:m1:exchange-traded-funds:premium}).
\end{example}

\begin{omfigure}
\begin{tikzpicture}
\begin{axis}[width=11cm,height=4.8cm,
  xlabel={time (arbitrary steps)},ylabel={premium (basis points)},
  tick label style={font=\small},label style={font=\small},
  xmin=0,xmax=120,ymin=-10,ymax=10,grid=major,grid style={gray!25},table/col sep=comma]
\addplot[omDef,thick,restrict x to domain=0:120] table[x=t,y=premium_bp]{figdata/markets-1/14-exchange-traded-funds/premium.csv};
\addplot[omThm,dashed,thick] coordinates {(0,7) (120,7)};
\addplot[omThm,dashed,thick] coordinates {(0,-7) (120,-7)};
\end{axis}
\end{tikzpicture}

\omcaption{A simulated \omterm{def:m1:exchange-traded-funds:nav}{premium}. Order flow pushes it around; whenever it
reaches the band of \cref{ex:m1:exchange-traded-funds:band} (dashed) an
\omterm{def:m1:exchange-traded-funds:ap}{authorised participant} creates or redeems and brings it halfway back. Inside
the band nothing forces it to zero. Data: the tutorial's
simulation.}\label{fig:m1:exchange-traded-funds:premium}
\end{omfigure}

Most ETF trading never reaches the primary market. A \omterm{def:m1:what-a-trading-firm-does:market-maker}{market maker} who sells
ETF shares hedges with futures, with the basket or with a correlated ETF, and
lets buyers and sellers net against each other; only the residual, often at
the end of the day, becomes a creation. Shares outstanding change little on a
day when turnover is many times the fund's size. Market making in ETFs, the
business of several of the firms in \cref{ch:m1:what-a-trading-firm-does}, is
treated in One Quant Book 11.

\section{When the inside cannot be priced}

The arbitrage argument needs a basket price. Two kinds of ETF do not have
one. For a fund of Japanese stocks trading in New York, the basket's market
is closed all day. For a fund of corporate bonds, most of the holdings have
not traded today at all: their ``prices'' are a vendor's estimates.

\begin{proposition}[A stale NAV produces a spurious premium]\label{prop:m1:exchange-traded-funds:stale}
Let the true value $V_t$ of the holdings move, and let the published NAV
adjust only partly each day, $N_t = N_{t-1} + (1-\theta)(V_t - N_{t-1})$ with
$0 < \theta < 1$. If the ETF trades at $V_t$, then after a sustained fall of
$V$ the ETF shows a discount to NAV, which closes without any trade in the
ETF being a bargain.
\end{proposition}
\begin{proof}
$V_t - N_t = \theta(V_t - N_{t-1})$. After a fall $V_t < N_{t-1}$, so $V_t <
N_t$: the ETF, at $V_t$, is below NAV. When $V$ stabilises the difference
decays geometrically at rate $\theta$.
\end{proof}

\begin{example}[March 2020]\label{ex:m1:exchange-traded-funds:march}
During the sell-off of March 2020, investment-grade corporate bond ETFs
closed about 365 basis points below their NAVs, on an asset-weighted basis,
on 12 March and again on 19 March; the largest of them went from a discount
of 2.8\% on 20 March to a \omterm{def:m1:exchange-traded-funds:nav}{premium} of 2.9\% on 23 March, the day a central-bank
purchase programme was announced. A central-bank study of the episode read
the discounts as the ETFs' prices leading stale bond valuations, in a market
where dealers had stopped making prices: the ETF was the price, and the NAV
the estimate.
\end{example}

\begin{omfigure}
\begin{tikzpicture}
\begin{axis}[width=11cm,height=5cm,
  xlabel={day},ylabel={value per share},
  tick label style={font=\small},label style={font=\small},
  xmin=0,xmax=59,ymin=84,ymax=102,grid=major,grid style={gray!25},
  legend columns=2,legend style={font=\footnotesize,at={(0.5,-0.3)},anchor=north,draw=none,fill=none,/tikz/every even column/.append style={column sep=8pt}},
  table/col sep=comma]
\addplot[omDef,very thick] table[x=day,y=true]{figdata/markets-1/14-exchange-traded-funds/stale.csv};
\addplot[omThm,very thick,dashed] table[x=day,y=nav]{figdata/markets-1/14-exchange-traded-funds/stale.csv};
\legend{ETF price (true value of the holdings),published NAV (stale valuations)}
\end{axis}
\end{tikzpicture}

\omcaption{A simulated ten-day sell-off in an illiquid asset class. The NAV
lags; the ETF appears to trade at a discount that reaches 4\% and then
``closes'' as the valuations catch up. Data: the tutorial's simulation, with
$\theta = 0.75$.}\label{fig:m1:exchange-traded-funds:stale}
\end{omfigure}

The practical rule: a \omterm{def:m1:exchange-traded-funds:nav}{premium} is evidence of mispricing only to the extent
that the basket can be traded at the prices used in the NAV. For domestic
equities it can, and seven basis points is the whole story. For bonds and
for closed markets the ETF is often the best available estimate of value, and
the right comparison is with futures, with other ETFs and with the few bonds
that did trade.

\section{Leveraged and inverse funds}

\begin{definition}[Leveraged ETF]\label{def:m1:exchange-traded-funds:leveraged}
A \emph{leveraged ETF}\index{leveraged ETF} with factor $\beta$ promises
$\beta$ times the \emph{daily} return of an index, before fees; an
\emph{inverse} ETF has $\beta < 0$. It obtains the exposure with swaps and
futures and resets it at each close.
\end{definition}

\begin{proposition}[The daily rebalance]\label{prop:m1:exchange-traded-funds:rebalance}
A fund with net assets $N$ and exposure $\beta N$ experiences an index return
$r$ during the day. To be $\beta$ times leveraged at the close it must
increase its exposure by
\[
\Delta = \beta(\beta - 1)\,N\,r .
\]
Since $\beta(\beta-1) > 0$ for every $\beta$ outside $[0,1]$, leveraged
\emph{and} inverse funds all buy after a rise and sell after a fall.
\end{proposition}
\begin{proof}
After the move the assets are $N(1+\beta r)$ and the exposure $\beta N(1+r)$.
The target exposure is $\beta N(1+\beta r)$. The difference is $\beta N(\beta
r - r)$.
\end{proof}

\begin{example}[Six percent of assets per percent]\label{ex:m1:exchange-traded-funds:six}
For $\beta = 3$, $\Delta = 6Nr$; for $\beta = -1$, $\Delta = 2Nr$; for $\beta
= -3$, $\Delta = 12Nr$. A \$10 billion three-times fund must buy \$3 billion of
exposure at the close of a day on which its index rose 5\%, whatever anyone
at the fund thinks. The flow is public arithmetic, is concentrated in the
final minutes, and has the sign of the day's move: it amplifies closing
moves, and anticipating it is one of the event strategies of One Quant Book 8
(\cref{fig:m1:exchange-traded-funds:rebalance}).
\end{example}

\begin{omfigure}
\begin{tikzpicture}
\begin{axis}[width=9.5cm,height=4.4cm,ybar,bar width=14pt,
  xlabel={leverage factor $\beta$},ylabel={trade (\% of net assets)},
  symbolic x coords={-3,-2,-1,1,2,3},xtick=data,
  tick label style={font=\small},label style={font=\small},
  ymin=0,ymax=14,ymajorgrids,grid style={gray!25},
  nodes near coords,nodes near coords style={font=\footnotesize},table/col sep=comma]
\addplot[fill=omDef!60,draw=omDef] table[x=beta,y=trade_pct_nav]{figdata/markets-1/14-exchange-traded-funds/rebalance.csv};
\end{axis}
\end{tikzpicture}

\omcaption{Exposure a fund must buy at the close after a $+1\%$ day, as a
percentage of its net assets: $\beta(\beta-1)$. Inverse funds trade more than
leveraged funds of the same magnitude, and in the same direction. Data:
\cref{prop:m1:exchange-traded-funds:rebalance}.}\label{fig:m1:exchange-traded-funds:rebalance}
\end{omfigure}

\begin{definition}[Volatility decay]\label{def:m1:exchange-traded-funds:decay}
\emph{Volatility decay}\index{volatility decay} is the shortfall of a
leveraged fund's multi-day return relative to $\beta$ times (or the $\beta$-th
power of) the index's return over the same period, caused by daily
resetting in a volatile market.
\end{definition}

\begin{proposition}[How much decay]\label{prop:m1:exchange-traded-funds:decay}
If the index follows $\dd S/S = \mu\dd t + \sigma\dd W$ and the fund is
continuously rebalanced to leverage $\beta$ with no financing cost, then
\[
\frac{L_T}{L_0} \;=\; \left(\frac{S_T}{S_0}\right)^{\beta}
\exp\!\Bigl(-\tfrac12\,\beta(\beta-1)\,\sigma^2 T\Bigr).
\]
\end{proposition}
\begin{proof}
$\dd L/L = \beta\,\dd S/S$, so by It\^o's formula $\dd\ln L = \beta\mu\dd t +
\beta\sigma\dd W - \tfrac12\beta^2\sigma^2\dd t$, while $\beta\,\dd\ln S =
\beta\mu\dd t + \beta\sigma\dd W - \tfrac12\beta\sigma^2\dd t$. Subtract and
integrate.
\end{proof}

With $\sigma = 30\%$ and $T = 1$ year the factor is $\exp(-0.27) = 0.76$ for
$\beta = 3$ and $\exp(-0.09) = 0.91$ for $\beta = -1$. The fund is not
cheating: it delivered $\beta$ times each day's return. It is the investor
who holds it for a year, expecting $\beta$ times the year's return, who has
misread the contract.

\begin{omfigure}
\begin{tikzpicture}
\begin{axis}[width=11cm,height=5cm,
  xlabel={trading day},ylabel={value of 1 invested},
  tick label style={font=\small},label style={font=\small},
  xmin=1,xmax=252,ymin=0.4,ymax=1.8,grid=major,grid style={gray!25},
  legend columns=3,legend style={font=\footnotesize,at={(0.5,-0.3)},anchor=north,draw=none,fill=none,/tikz/every even column/.append style={column sep=8pt}},
  table/col sep=comma]
\addplot[omExo,very thick] table[x=day,y=index]{figdata/markets-1/14-exchange-traded-funds/leveraged.csv};
\addplot[omThm,very thick] table[x=day,y=three_x]{figdata/markets-1/14-exchange-traded-funds/leveraged.csv};
\addplot[omDef,thick,dashed] table[x=day,y=naive_three_x]{figdata/markets-1/14-exchange-traded-funds/leveraged.csv};
\legend{index,three-times daily fund,three times the index's return}
\end{axis}
\end{tikzpicture}

\omcaption{One simulated year with a daily volatility of 2\% and daily returns
that sum to zero. The index ends 5.8\% lower; three times that would be
$-17.4\%$; the three-times fund ends 41.8\% lower. Data: the tutorial's
simulation.}\label{fig:m1:exchange-traded-funds:decay}
\end{omfigure}

\section{Tutorial: bands, stale values and decay}\label{tut:m1:exchange-traded-funds:bands}

\begin{tutorial}
\textbf{Goal.} Compute an arbitrage band, reproduce a spurious discount from
a stale NAV, and measure \omterm{def:m1:exchange-traded-funds:decay}{volatility decay} against the formula. \textbf{End
state:} the four data figures of this chapter.
\begin{tutsteps}
\item \textbf{The band} is a sum of four costs.
\omcode{code/markets-1/14-exchange-traded-funds/python/etf_demo.py}{9}{19}{The costs an authorised participant must cover before an arbitrage pays.}\label{lst:m1:exchange-traded-funds:band}
\item \textbf{Stale NAV.} One line of smoothing is enough to manufacture a
4\% discount in a ten-day sell-off.
\omcode{code/markets-1/14-exchange-traded-funds/python/etf_demo.py}{39}{45}{A net asset value computed from quotes that adjust only partly each day.}\label{lst:m1:exchange-traded-funds:stale}
\item \textbf{The rebalance and the decay.}
\omcode{code/markets-1/14-exchange-traded-funds/python/etf_demo.py}{48}{60}{The closing trade of a leveraged fund, its path, and the continuous-time decay factor.}\label{lst:m1:exchange-traded-funds:lev}
\item \textbf{Check the formula.} The tests simulate 4\,000 years at 30\%
volatility and compare the mean log ratio of the three-times fund to the
cube of the index with $-\tfrac12\cdot6\cdot0.09 = -0.27$.
\end{tutsteps}
\textbf{What to change next.} Add a financing spread and a 0.95\% annual fee
to the leveraged fund. Then make the index trend (positive autocorrelation of
daily returns) and find the regime in which the three-times fund
\emph{beats} three times the index.
\end{tutorial}

\section{Build: the premium monitor}\label{bld:m1:exchange-traded-funds:monitor}

\begin{build}
\bfield{Purpose} The miniature firm makes markets in ETFs on its simulated
exchange (One Quant Book 11). It needs a live fair value for each fund and an
alarm when the market price leaves the arbitrage band.
\bfield{Interface} \texttt{Basket(etf, shares\_per\_unit, components, cash)}
with components as (symbol, quantity per \omterm{def:m1:exchange-traded-funds:ap}{creation unit});
\texttt{indicative\_nav(basket, mids)}; \texttt{premium\_bp(price, inav)};
\texttt{band\_bp(costs)}; \texttt{signal(price, inav, costs)} returning
\texttt{create}, \texttt{redeem} or \texttt{none}; and for each component a
\texttt{stale} flag when its last quote is older than a threshold.
\bfield{Rules} Integer prices; the indicative value uses mids from the
consolidated quote of \cref{bld:m1:us-equity-market-structure:nbbo}. If
components worth more than a configurable fraction of the basket are stale
the monitor reports \texttt{unreliable} instead of a \omterm{def:m1:exchange-traded-funds:nav}{premium}. Corporate
actions change the basket overnight
(\cref{bld:m1:shares-corporate-actions-indices:adjuster}).
\bfield{Acceptance tests} \texttt{code/firm/etfmonitor/tests/}: a three-stock
basket priced by hand; signals at and beyond the band; the stale-component
rule.
\bfield{Stretch} Replace components whose market is closed by a regression on
instruments that are open (futures, currency, a related ETF): the fair value
of an international fund.
\end{build}

\begin{omsources}
\item ETFGI, press releases on global ETF industry assets, July--September
2026 (as reported by Markets Media, ``Global ETF Assets Top \$24 Trillion for
First Time'').
\item US Securities and Exchange Commission, \emph{SEC Adopts New Rule to
Modernize Regulation of Exchange-Traded Funds}, press release 2019-190, September 2019.
\item S.~Aramonte and F.~Avalos, ``The recent distress in corporate bond
markets: cues from ETFs'', \emph{BIS Bulletin} no.~6, 14 April 2020; K.~Todorov,
``The anatomy of bond ETF arbitrage'', \emph{BIS Quarterly Review}, March
2021.
\item M.~Cheng and A.~Madhavan, ``The dynamics of leveraged and inverse
exchange-traded funds'', \emph{Journal of Investment Management} 7 (2009).
\item A.~Madhavan, \emph{Exchange-Traded Funds and the New Dynamics of
Investing}, Oxford University Press, 2016.
\end{omsources}

\section{Exercises}

\begin{exercise}[$\star$]\label{exo:m1:exchange-traded-funds:1}
An ETF's NAV is \$82.40 and it trades at \$82.31. Give the \omterm{def:m1:exchange-traded-funds:nav}{premium} in basis
points. With a band of 9 basis points, does an \omterm{def:m1:exchange-traded-funds:ap}{authorised participant} act,
and how?
\end{exercise}

\begin{exercise}[$\star$]\label{exo:m1:exchange-traded-funds:2}
A \omterm{def:m1:exchange-traded-funds:ap}{creation unit} is 50\,000 ETF shares at a NAV of \$40, and the fund charges
\$500 per creation. Express the fee in basis points of the unit's value. What
is it if the AP creates ten units in one order for the same fee?
\end{exercise}

\begin{exercise}[$\star$]\label{exo:m1:exchange-traded-funds:3}
A $-2\times$ fund has net assets of \$800 million and its index falls 4\%
today. Give the fund's return, its new net assets, and the trade it must make
at the close, with its direction.
\end{exercise}

\begin{exercise}[$\star\star$]\label{exo:m1:exchange-traded-funds:4}
An index rises 10\% on day one and falls 9.09\% on day two, ending where it
started. Compute the two-day return of a $2\times$ fund, a $3\times$ fund and
a $-1\times$ fund. Which loses most, and why?
\end{exercise}

\begin{exercise}[$\star\star$]\label{exo:m1:exchange-traded-funds:5}
For $\sigma = 40\%$ a year compute the decay factor of
\cref{prop:m1:exchange-traded-funds:decay} over one year for $\beta = 2$, $3$
and $-2$. Over what horizon does a $3\times$ fund lose 10\% to decay at that
volatility?
\end{exercise}

\begin{exercise}[$\star\star$]\label{exo:m1:exchange-traded-funds:6}
A bond ETF's NAV follows \cref{prop:m1:exchange-traded-funds:stale} with
$\theta = 0.8$. The true value falls from 100 to 94 in one day and stays
there. Give the NAV and the apparent discount on that day and on the next
three days. After how many days is the discount below 0.5\%?
\end{exercise}

\begin{exercise}[$\star\star\star$]\label{exo:m1:exchange-traded-funds:7}
\emph{Coding.} With \texttt{simulate\_premium} and the costs of
\cref{ex:m1:exchange-traded-funds:band}, run 5\,000 steps with seed 1 and
report the number of steps at which an \omterm{def:m1:exchange-traded-funds:ap}{authorised participant} acted and the
mean absolute \omterm{def:m1:exchange-traded-funds:nav}{premium}. Halve the basket cost $c_B$ and report the same. What
does an investor in the ETF gain?
\end{exercise}

\begin{exercise}[$\star\star\star$]\label{exo:m1:exchange-traded-funds:8}
\emph{Find the flaw.} A newsletter writes: ``This high-yield bond ETF closed
at a 3\% discount to NAV in a falling market. Buy it: the discount always
closes within two weeks, a riskless 3\%.'' Explain why the discount closing
does not imply a profit, and describe the position that would be needed to
capture it if it were real.
\end{exercise}

\section{Problem: The Daily Rebalance of a Three-Times Fund}

\begin{problem}[{Weekend problem --- the last ten minutes of a leveraged fund}]\label{pb:m1:exchange-traded-funds:1}
A $3\times$ fund and a $-3\times$ fund track the same equity index. At last
night's close the first had net assets of \$8 billion and the second \$3
billion. The index futures trade \$300 billion a day, of which a tenth in the
last ten minutes.

\medskip\noindent\textbf{Part I --- A quiet day.}
\begin{enumerate}
\item The index rises 1\%. Give each fund's return and new net assets.
\item Give each fund's closing trade with its direction, and the total.
\item Express the total as a fraction of the last ten minutes' volume.
\item With the square-root rule of \cref{met:m1:the-sell-side:sqrt}, a daily
volatility of 1\% and $\kappa = 0.7$, estimate the impact of that trade,
taking $V$ to be the ten-minute volume and scaling $\sigma$ to ten minutes of
a 390-minute day.
\end{enumerate}

\medskip\noindent\textbf{Part II --- A violent day.}
\begin{enumerate}[resume]
\item The index falls 5\%. Give each fund's return and new net assets.
\item Give the closing trades and the total.
\item As a fraction of the last ten minutes' volume?
\item Estimate the impact as in question 4 with a daily volatility of 3\%.
\item The impact feeds back: if the funds' selling moves the index a further
$x$, they must sell a little more. Using your answer to question 8 as $x$,
by how much does the required trade grow?
\end{enumerate}

\medskip\noindent\textbf{Part III --- A bad week.}
The index returns $-5\%$, $+4\%$, $-6\%$, $+5\%$, $+2.5\%$ over five days.
\begin{enumerate}[resume]
\item Give the index's cumulative return.
\item Give the $3\times$ fund's cumulative return, and three times the
index's.
\item Give the $-3\times$ fund's cumulative return, and minus three times the
index's.
\item An investor held both funds in equal dollar amounts ``to be hedged''.
Give her return.
\item Explain her result with \cref{prop:m1:exchange-traded-funds:decay}.
\end{enumerate}

\medskip\noindent\textbf{Part IV --- Judgement.}
\begin{enumerate}[resume]
\item Who is on the other side of the closing trades of Part II, and what do
they need to know in advance?
\item Why does the $-3\times$ fund trade twice as much per dollar of assets
as the $3\times$ fund?
\item What daily index move would wipe out the $3\times$ fund? How do such
funds' prospectuses deal with it?
\item A regulator asks whether these funds destabilise the close. Give the
quantity that answers the question and its value on the violent day.
\item State the \emph{named result}: the notional the two funds together must
trade at the close after the 5\% fall, in billions of dollars.
\item In one sentence: for whom is a \omterm{def:m1:exchange-traded-funds:leveraged}{leveraged ETF} a sensible instrument?
\end{enumerate}
\end{problem}

\section{Interview questions}

\begin{interviewq}[$\star$ \iqroles{trader, researcher, developer}]\label{iq:m1:exchange-traded-funds:1}
How does an ETF's price stay close to the value of what it holds?
\end{interviewq}

\begin{interviewq}[$\star$ \iqroles{trader, researcher}]\label{iq:m1:exchange-traded-funds:2}
A $2\times$ \omterm{def:m1:exchange-traded-funds:leveraged}{leveraged ETF}'s index is up 10\% over the year. Is the ETF up
20\%?
\end{interviewq}

\begin{interviewq}[$\star\star$ \iqroles{trader, researcher}]\label{iq:m1:exchange-traded-funds:3}
An ETF holding Japanese stocks trades in New York at a 1.5\% \omterm{def:m1:exchange-traded-funds:nav}{premium} to its
NAV at 11:00 New York time. Is that an arbitrage?
\end{interviewq}

\begin{interviewq}[$\star\star$ \iqroles{trader, researcher}]\label{iq:m1:exchange-traded-funds:4}
You make markets in an ETF of 500 stocks. A client buys \$50 million from
you. Walk me through your hedge and your end of day.
\end{interviewq}

\begin{interviewq}[$\star\star$ \iqroles{researcher, mle}]\label{iq:m1:exchange-traded-funds:5}
Derive the amount a leveraged fund must trade at the close as a function of
its leverage and the day's return. What is surprising about inverse funds?
\end{interviewq}

\begin{interviewq}[$\star\star\star$ \iqroles{researcher, trader}]\label{iq:m1:exchange-traded-funds:6}
In a crisis a bond ETF trades 4\% below NAV. Build the case that the ETF is
mispriced, the case that the NAV is, and say what data would decide.
\end{interviewq}
